\section{Background and Motivation}
\label{sec:background}
\subsection{The Observation--Action Loop}
A VLA policy consumes images, a language instruction, and any required
robot state, and returns actions. In a policy that uses action chunks,
one request produces several future control steps; OpenVLA-OFT explicitly
uses action chunking in its fine-tuning recipe~\cite{openvlaoft}. The
actuator control rate and the policy request rate therefore describe
different parts of the system. Increasing the former does not imply
that a fresh visual inference is performed at every control step.

This distinction matters when evaluating an edge--cloud path. An upload
can complete faster without improving a task if the returned action is
not needed immediately, or if errors in the visual representation impair
the next motion. Conversely, a delayed request can make its observation
stale by the time the action is used. TimelyLLM studies the relationship
between inference service and physical execution for agents~\cite{timelyllm}.
Here we focus on the representation sent with visual policy requests.
Our measured request latency ends when the action response returns; it
does not by itself measure observation age at physical action application.

\subsection{The Cost of Visual Communication}
\label{sec:bg-cost}
There are three relevant paths. Image offloading encodes pixels on the
edge and reconstructs them before cloud visual inference. Conventional
partitioning runs a model prefix locally and transmits intermediate
activations. Our path learns a compact representation and a receiver
that recovers policy input tokens. Table~\ref{tab:paths} summarizes these
interfaces. Model splitting and feature compression are established
ideas~\cite{neurosurgeon,bottlenetpp}; the question here is how to realize
the feature interface for a frozen VLA on robot hardware.

\begin{table}[t]
\caption{Visual communication interfaces. The policy after visual-token
construction is shared within each backbone comparison.}
\label{tab:paths}
\centering\small
\begin{tabularx}{\columnwidth}{lXX}
\toprule
Path & Transmitted representation & Cloud visual work \\
\midrule
Image codec & Image/video bitstream & Image decoding and original visual encoder \\
Model split & Intermediate activations & Remaining model computation \\
Feature streaming & Quantized feature bitstream & Feature recovery and token injection \\
\bottomrule
\end{tabularx}
\end{table}

The relevant cost includes edge preparation, transfer, and cloud
processing. A method with fewer bytes can still have a longer request
latency if encoding is expensive. A method with a more expensive encoder
can be useful if it removes sufficient receiver work. These alternatives
must be timed on the intended hardware. In particular, savings measured
on one cloud GPU need not persist on another.

\subsection{Why Policy Alignment Matters}
\label{sec:bg-align}
Image compression commonly trades rate against visual distortion;
learned hyperprior codecs provide one example of this
formulation~\cite{hyperprior}. A VLA introduces a different endpoint:
the recovered observation must support the action policy. SPARC addresses
this distinction with action distortion and a rate penalty
rather than relying on image fidelity alone~\cite{sparc}. Our design
retains the task-oriented endpoint while changing the receiver output
from pixels to policy tokens.

Our diagnostics also show that additional capacity or lower
feature error cannot be interpreted in isolation. Changing an input
preprocessing protocol changes the policy's reference performance, and
perturbed camera views can expose a quality loss that is small on the
original benchmark. We therefore report the input protocol, retain
negative results, and measure closed-loop simulation success separately
from feature reconstruction error and request latency.
